\documentclass[letterpaper,11pt]{article}

\usepackage[empty]{fullpage}
\usepackage[hidelinks]{hyperref}
\usepackage{fontawesome5}
\usepackage{enumitem}
\usepackage{parskip}

% Adjust margins
\addtolength{\oddsidemargin}{-0.4in}
\addtolength{\evensidemargin}{-0.4in}
\addtolength{\textwidth}{0.8in}
\addtolength{\topmargin}{-0.5in}
\addtolength{\textheight}{1.0in}

\ifdefined\pdfgentounicode
  \input{glyphtounicode}
  \pdfgentounicode=1
\fi

\begin{document}

\begin{center}
    {\Huge \scshape Juan Becerra} \\ \vspace{4pt}
    \textbf{\large Software Engineer -- Gameplay, Animation \& Tools} \\ \vspace{4pt}
    \small Seattle, WA ~ $|$ ~
    \raisebox{-0.1\height}\faPhone\ (208) 550-9191 ~ $|$ ~
    \href{mailto:becerrajuan007@gmail.com}{\raisebox{-0.2\height}\faEnvelope\ \underline{becerrajuan007@gmail.com}} ~ $|$ ~
    \href{https://hoowan.dev}{\raisebox{-0.2\height}\faGlobe\ \underline{hoowan.dev}} \\ \vspace{2pt}
    \href{https://linkedin.com/in/hoowan/}{\raisebox{-0.2\height}\faLinkedin\ \underline{linkedin.com/in/hoowan}} ~ $|$ ~
    \href{https://github.com/hoowan-dev}{\raisebox{-0.2\height}\faGithub\ \underline{github.com/hoowan-dev}}
\end{center}

\vspace{4pt}
\hrule
\vspace{14pt}

\noindent October 26, 2026

\vspace{8pt}
\noindent \textbf{Samy Duc} \\
EA Studios - Respawn \\
Electronic Arts \\
Los Angeles, CA / Vancouver, BC

\vspace{8pt}
\noindent \textbf{Re: Gameplay Data Science Engineer (Apex Legends) (Role ID: 216112)}

\vspace{8pt}
\noindent Dear Samy Duc,

I am writing to express my strong interest in the \textbf{Gameplay Data Science Engineer} 
position at Respawn Entertainment. Over the past five years at Electronic Arts (four years as a 
Gameplay \& Animation Software Engineer at \textbf{Maxis} on \emph{Project Rene} (\emph{The Sims}) 
and currently at \textbf{Lifestyle}), combined with my ongoing graduate coursework in Machine 
Learning \& Deep Learning at the University of Washington, my work has centered on C++ 
engine-level gameplay systems, agent navigation and behavior architectures, and Python/PyTorch 
machine learning pipelines. I would be thrilled to help advance intelligent bot behaviors and 
scalable AI frameworks for Apex Legends.

In my current role at \textbf{EA Lifestyle}, I bridge runtime engine data extraction and machine 
learning workflows, building a high-throughput Unreal Engine C++ plugin that captures, structures, 
and exports runtime character trajectory, pose, and velocity datasets to train and populate 
motion-matching query databases. In parallel, I prototyped and integrated ML-assisted pipelines 
using Python, PyTorch, and ONNX Runtime into authoring workflows for automated asset tagging and 
procedural pose generation, and architected cross-studio Python data-processing tools that reduced 
manual workflows by over 70\%.

During my tenure at \textbf{EA Maxis Studios}, I focused on engine-level agent navigation, 
locomotion, and behavior tooling in Unreal Engine 5 and C++. I architected a multi-threaded, 
asynchronous path-planning and navigation API that enabled deterministic agent routing and 
real-time obstacle avoidance for shared multiplayer simulation. I also implemented core locomotion 
and runtime animation systems, including motion-matching, direct/indirect control schemes, and 
custom IK solvers, and partnered closely with gameplay designers to build in-editor tuning tools 
and visual debuggers while profiling runtime memory and compute across character evaluation graphs 
and spatial queries.

As an internal EA engineer based in the Pacific Northwest with a deep appreciation for Apex 
Legends' movement and combat depth, I am excited about the opportunity to collaborate with 
Respawn's team to translate creative design goals into high-performance, realistic agent behaviors. 
Thank you for your time and consideration, and I look forward to discussing how my experience can 
contribute to the team.

\vspace{10pt}
\noindent Sincerely,

\vspace{14pt}
\noindent \textbf{Juan Becerra}

\end{document}
